% ============================================================
%  APPENDIX
%  Supplementary material organized by chapter/paper.
% ============================================================

\appendix

% ---- Mango Supplementary Material ----
\chapter{Mango Supplementary Material}
\label{app:mango}

\section{Additional Firmware Tests}
\label{app:mango-additional-firmware}

To demonstrate extensibility beyond the \karonteDataset, we randomly selected six firmware samples from different vendors (shown in Table~\ref{tab:appendix_firmware_test}).
No additional effort or analyses were added to \toolname to generate these results; they were produced in the same manner as the results discussed in the evaluation (Section~\ref{sec:evaluation}).
\toolname found a significant number of \trupoc{}s in these firmware samples, including \trupoc{}s in Belkin and Linksys firmware images that were not tested in the \karonteDataset.
Each generated \trupoc was manually analyzed to determine the overall true positive rate, which was 41.3\%.

\begin{table}[H]
    \centering
    \small
    \begin{tabular}{lr:rrrrr}
        \hline
        \textbf{Vendor} & \textbf{Device} & \multicolumn{5}{:c}{\textbf{\toolname}}                                                                             \\
                        &                 & \textbf{\trupoc}                        & \textbf{TP} & \textbf{Time} & \textbf{\trupoc Bins} & \textbf{Total Bins} \\
        \hline
        D-Link          & DIR600          &
        1               & 0               & 0:47h                                   & 1           & 48                                                          \\

        D-Link          & DIR300          &
        1               & 0               & 0:44h                                   & 1           & 40                                                          \\

        Belkin          & F9J1102         &
        52              & 25              & 1:46h                                   & 4           & 71                                                          \\

        Linksys         & WRT320N         &
        22              & 3               & 0:07h                                   & 4           & 49                                                          \\

        Trendnet        & TEW733          &
        4               & 0               & 0:50h                                   & 2           & 58                                                          \\

        Tenda           & N60             &
        109             & 50              & 3:25h                                   & 8           & 65                                                          \\
        \hline
        Total           & -               &
        189             & 78              & 7:39h                                   & 20          & 331                                                         \\
        \hline
    \end{tabular}
    \caption{
        \toolname analysis on six firmware images outside the dataset of \karonte and \satc.
    }
    \label{tab:appendix_firmware_test}
\end{table}

% ---- Artiphishell Supplementary Material ----
\chapter{Artiphishell Supplementary Material}
\label{app:artiphishell}

\section{Services}
\label{app:artiphishell-services}

The different components of \artiphishell{} rely on various services to operate.
These services are initialized and dynamically configured once the required pre-processing stages have completed and the necessary artifacts are available.

\subsection{Function Resolver}
\label{app:function-resolver}

Locating the original source-code implementations of functions is a critical task for all components in \artiphishell{}.
While this task appears rudimentary at first glance, real code-bases commonly include build-steps that complicate a direct mapping of compilation-artifact function representations back to the original code.

Notable examples of this include auto-generated code (e.g., in parser-generators like Flex, YACC, etc.), build and link-time performance optimizations (e.g., SQLite's ``amalgamation'') and more.

\section{Nautilus Grammar Formats}
\label{app:nautilus-grammar-formats}

\begin{figure}[H]
\centering
\begin{lstlisting}[language={},basicstyle=\scriptsize\ttfamily,frame=single,breaklines=true]
{ "rules": {
    "START": [[{"NT": "HTML"}]],
    "HTML":  [[{"T": "<html>"}, {"NT": "HEAD"}, {"NT": "BODY"}, {"T": "</html>"}]],
    "HEAD":  [[{"T": "<head>"}, {"NT": "TITLE"}, {"T": "</head>"}]],
    "TITLE": [[{"T": "<title>"}, {"NT": "TEXT"}, {"T": "</title>"}]],
    "BODY":  [[{"T": "<body>"}, {"NT": "CONTENT"}, {"T": "</body>"}]],
    "CONTENT": [
      [{"NT": "PARAGRAPH"}],
      [{"NT": "HEADING"}],
      [{"NT": "PARAGRAPH"}, {"NT": "CONTENT"}]],
    "HEADING":   [[{"T": "<h1>"}, {"NT": "TEXT"}, {"T": "</h1>"}]],
    "PARAGRAPH": [[{"T": "<p>"}, {"NT": "TEXT"}, {"T": "</p>"}]],
    "TEXT": [[{"T": "Hello World"}]]
}}
\end{lstlisting}
\caption{Nautilus JSON grammar format}
\label{fig:grammar-format-json}
\end{figure}

\begin{figure}[H]
\centering
\begin{lstlisting}[language=Python,basicstyle=\scriptsize\ttfamily,frame=single]
# Basic HTML Grammar
ctx.rule("START", b"{HTML}")

# HTML structure
ctx.rule("HTML", b"<html>{HEAD}{BODY}</html>")

# Head section
ctx.rule("HEAD", b"<head>{TITLE}</head>")
ctx.rule("TITLE", b"<title>{TEXT}</title>")

# Body section
ctx.rule("BODY", b"<body>{CONTENT}</body>")

# Content alternatives
ctx.rule("CONTENT", b"{PARAGRAPH}")
ctx.rule("CONTENT", b"{HEADING}")
ctx.rule("CONTENT", b"{PARAGRAPH}{CONTENT}")

# Elements
ctx.rule("HEADING", b"<h1>{TEXT}</h1>")
ctx.rule("PARAGRAPH", b"<p>{TEXT}</p>")

# Text content
ctx.literal("TEXT", b"Hello World")
\end{lstlisting}
\caption{Nautilus Python grammar format}
\label{fig:grammar-format-python}
\end{figure}
